\section{Four campaign case studies}
\label{sec:cases}

\subsection{The decoy-gradient experiment}
\label{sec:f1}

The decoy-gradient experiment mixes real activation rows with decoys.
The trusted client sends protected activations to the cloud and receives
its output. The client holds the private loss and returns the output
gradient to the cloud. Decoys receive exactly zero gradient because the
loss ignores them. That support pattern reveals which rows bear the loss,
so the disclosure concerns the real/decoy partition. Token identities
require a separate content attack~\cite{paper1} (\artifact{E1}).

The experiment records report exact partition recovery on 4,096 of 4,096
inspected frames \emph{per run} across nine seeds. A frequent token content
probe separately improves on a constant guess baseline by $+0.65$ to
$+1.50$ percentage points. The two results use different denominators.
In six cells that satisfy the utility rule and expose gradients, three
seeds on each of two datasets pass the privacy gate for forward
observations alone. The joint view exceeds the privacy gate of $+1.0$
percentage points set in advance (\artifact{E1}).

The result for forward observations remains valid within that view.
Extending its non-detection to the cloud's joint view exceeds the scope of
the test, where the additional observation exposes a breach. The records
also report that clipping each gradient row and adding Gaussian noise
removes the scored partition disclosure. The partition is disclosed in
0 of 1,024 inspected frames, with agreement 0.400 at a 32/80 base rate.
The associated held out cross entropy \emph{utility cost} is approximately
$0.01$ nats. This quantity measures utility rather than residual leakage.
The mitigation result covers the scored probes, with attacks accumulating
observations across training and other omitted families still untested.
It nevertheless supplies evidence against dismissing all perturbation
based mitigation.

\subsection{The boundary-replacement audit}
\label{sec:f2}

The boundary-replacement audit found that representations of
$(x_1,\ldots,x_{T-1})$ were scored against labels
$(x_2,\ldots,x_T)$.
A row representing $x_t$ was therefore evaluated against $x_{t+1}$.
The scores measure predictive information about the next token, while
the privacy claim concerned reconstruction of the represented input.
Even a high positive control score under this alignment validates only
the predictive target. The scores can be retained as next token probe
results, but their use to establish current token privacy is invalid
(\artifact{B1}).

Earlier independent audit notes report that separately trained controls
with a zero return match or outperform the trained remote model in all
six clean comparisons pairing corpus and seed. We retain this result as
\emph{source-reported}. Learned remote computation was therefore not shown
to be necessary for the observed utility endpoints. Removing the return
from an already trained model is a different intervention, and can affect
that model even when separate training without a return meets the utility
target (\artifact{B1}).

One matched checkpoint evaluation also exposes a difference between
utility metrics. NLL rises by approximately $0.034$ nats and passes its
rule, while teacher forced answer token top-1 match falls from $34.79\%$
to $18.15\%$. The drop is $16.64$ percentage points. This specific measure
of answer token matching differs from standard multiple choice task
accuracy, and the observation comes from one checkpoint rather than
replicated evaluation of the deployed protocol. The records therefore
narrow the quality and delegation claims while leaving corrected
reconstruction results unavailable (\artifact{B1}).

\subsection{The privacy, utility and gain exploration}
\label{sec:f3}

The privacy, utility and gain exploration completed 232 runs and rejected
all 19 configurations under its frozen exploratory rules. Reconciliation
of the cell table yields 13 qualified privacy breaches, five utility gate
failures and seven learning gain failures. The rejection reasons overlap
across configurations. Six configurations have undetermined privacy and
independent rejection on utility or gain, with width-64 rejected on gain
alone (\artifact{R1}).

Four of the 24 required instruments qualified. These were input byte
probes on inbound and outbound gradients in the two recorded defence
categories. Missing coverage prevents a joint scoped privacy pass, but
the applicable qualified probes still establish breaches. Independent
utility and gain failures also retain their force. The records thus
describe an empty feasible set for the finite tested family under its
recorded rules, alongside incomplete privacy coverage (\artifact{R1}).

The span-6 cell shows why rejection reasons must be read with the exact
inequality. Its mean utility cost across four seeds is $0.0999705344$
nats, with an upper 95\% Student $t$ limit of $0.1195024415$.
Both pass their respective thresholds of $0.10$ and $0.12$. The mean is
$2.95\times10^{-5}$ nats below its threshold. Gain failure and a qualified
privacy breach account for the rejection. These intervals describe only
four seeds, without adjustment for multiple comparisons. The small cohort
limits the conclusion to the tested family (\artifact{R1}).

The proposed confirmation packet called for eight seeds. It was never
frozen, admitted or executed. Its power upper bound of $0.263772$ for the
whole family is conditional on exploratory means and standard deviations,
with substantial sensitivity to the variance assumptions. That calculation
served planning. It supplies no probability that the 19 rejections are
false and no confirmatory outcome evidence (\artifact{R1}).

\subsection{The defence-redesign pilot}
\label{sec:f4}

The defence-redesign pilot completed 66 runs under a frozen exploratory
protocol. It covered seven cells with three paired seeds per cell and
panels of 16,384 windows. No cell met the eligibility rules. All seven
pass the remote gain rule and five pass utility. The one-bit and two-bit
quantisation configurations fail utility. The composite privacy verdicts record four failed cells
and three undetermined cells. These composite verdicts retain the original
qualification policy (\artifact{P4}).

Qualified probes still provide substantial negative evidence about the
tested defences. Reported punctuation property advantages for the worst
seed range from $0.79$ to $0.89$ where qualified, against acceptance at
$0.02$. Advantages for the next byte class range from $0.097$ to $0.112$.
Punctuation uses correlation advantage, while next byte structure uses
classification accuracy advantage. The next byte is an intentional target
here. In the boundary-replacement audit, the defect arose from using it
to claim resistance to current token reconstruction. The large measured
signals support rejection of these defence settings. The ratio between a
signal and the acceptance threshold does not diagnose instrument failure
or an unattainable gate (\artifact{P4}).

The calibration constants serve distinct purposes. The positive control
floor is $0.10$ for properties and $0.02$ for structure. The null band
is $0.02$ and privacy acceptance is $0.02$. A positive control floor does
not bound undetected leakage from above. The qualification records treat
marginal null excursions as evidence about reliability, with no finding
of instrument breakage. The finite negative result therefore concerns
the tested mechanisms, strengths, budgets and qualified targets. Untested
attacks, general token privacy and universal impossibility remain outside
that conclusion (\artifact{P4}).
